\chapter{Clase 4: Covarianza y correlación}

\section{Covarianza entre dos variables}

Consideremos dos variables cuantitativas $X$ y $Y$, observadas conjuntamente en $n$ unidades. Para cada observación $i$, el par $(x_i,y_i)$ indica los valores de ambas variables en la misma unidad. La covarianza describe si, en promedio, las desviaciones de $X$ y de $Y$ respecto de sus medias tienden a tener el mismo signo o signos opuestos.

\subsection{Definición poblacional}

Si la población contiene $N$ pares $(x_i,y_i)$, su covarianza se define como

$$
\operatorname{Cov}(X,Y) = \sigma_{XY} = \frac{1}{N}\sum_{i=1}^{N}(x_i-\mu_X)(y_i-\mu_Y),
$$

donde $\mu_X$ y $\mu_Y$ son las medias poblacionales de $X$ y $Y$. La sumatoria agrega, para cada individuo, el producto de las dos desviaciones con respecto a sus medias.

\subsection{Covarianza muestral}

Cuando se dispone de una muestra de $n$ pares, se utiliza la covarianza muestral

$$
s_{XY} = \operatorname{cov}(X,Y) = \frac{1}{n-1}\sum_{i=1}^{n}(x_i-\bar{x})(y_i-\bar{y}),
$$

donde $\bar{x}$ y $\bar{y}$ son las medias muestrales. El divisor $n-1$ es el habitual para estimar la covarianza poblacional a partir de una muestra.

\section{El signo de la covarianza en un diagrama de dispersión}

En un diagrama de dispersión, las líneas $x=\bar{x}$ y $y=\bar{y}$ separan el plano en cuatro regiones. El signo de la contribución de cada observación a la suma de covarianza depende del producto $(x_i-\bar{x})(y_i-\bar{y})$:

\begin{itemize}
    \item Si $x_i>\bar{x}$ y $y_i>\bar{y}$, ambas desviaciones son positivas y su producto es positivo.
    \item Si $x_i<\bar{x}$ y $y_i<\bar{y}$, ambas desviaciones son negativas y su producto también es positivo.
    \item Si una observación queda por encima de una media y por debajo de la otra, las desviaciones tienen signos opuestos y su producto es negativo.
\end{itemize}

La figura ilustra cómo se suman estas contribuciones. Una covarianza positiva aparece cuando predominan los productos positivos; una negativa, cuando predominan los productos negativos. Los puntos exactamente sobre una de las medias aportan cero.

\begin{figure}[htbp]
    \centering
    \begin{tikzpicture}[x=1.25cm,y=1.05cm,every node/.style={font=\small}]
    % Ejes y medias de la nube de datos: \bar{x}=5 y \bar{y}=4.5.
    \draw[->,thick] (0,0) -- (10.6,0) node[right] {$x$};
    \draw[->,thick] (0,0) -- (0,9.1) node[above] {$y$};

    \foreach \x in {1,...,10}
        \draw (\x,0.08) -- (\x,-0.08) node[below=3pt] {\x};
    \foreach \y in {1,...,8}
        \draw (0.08,\y) -- (-0.08,\y) node[left=3pt] {\y};

    \draw[densely dashed,gray!75,thick] (5,0) -- (5,8.5)
        node[above,black] {$x=\bar{x}$};
    \draw[densely dashed,gray!75,thick] (0,4.5) -- (10,4.5)
        node[right,black] {$y=\bar{y}$};

    % Puntos simétricos respecto de (5,4.5), de modo que las medias son exactas.
    \foreach \x/\y in {
        1.5/1.7, 8.5/7.3, 2.2/3.0, 7.8/6.0, 2.9/2.5, 7.1/6.5,
        3.8/4.0, 6.2/5.0, 4.4/3.6, 5.6/5.4, 5.8/4.2, 4.2/4.8,
        6.7/4.9, 3.3/4.1, 7.4/5.8, 2.6/3.2, 8.3/6.6, 1.7/2.4
    }
        \fill[blue!70!black] (\x,\y) circle (2.7pt);

    % Observación resaltada en el cuadrante superior derecho.
    % Proyecciones de la observacion resaltada sobre los ejes.
    \draw[dotted,orange!80!black,thick] (7.4,0) -- (7.4,5.8);
    \draw[dotted,orange!80!black,thick] (0,5.8) -- (7.4,5.8);
    \fill[orange!80!black] (7.4,0) circle (1.5pt);
    \fill[orange!80!black] (0,5.8) circle (1.5pt);
    \node[below=15pt,text=orange!60!black] at (7.4,0) {$x_i$};
    \node[left=8pt,text=orange!60!black] at (0,5.8) {$y_i$};
    \draw[orange!85!black,thick] (7.4,5.8) circle (5pt);
    \draw[->,orange!85!black,thick] (7.55,5.95) -- (7.9,7.8);
    \node[align=left,anchor=west,text=orange!60!black] at (8.1,8.35)
        {$x_i>\bar{x},\quad y_i>\bar{y}$\\
         $(x_i-\bar{x})(y_i-\bar{y})>0$};
\end{tikzpicture}
    \caption{Las desviaciones respecto de las medias tienen productos positivos en los cuadrantes superior derecho e inferior izquierdo, y negativos en los otros dos.}
    \label{fig:signo-covarianza}
\end{figure}

La covarianza informa la dirección de la variación conjunta, pero su magnitud depende de las unidades de medición. Por ello, no existe un límite fijo común para sus valores y no debe compararse directamente entre pares de variables medidos en escalas distintas.

\section{Correlación lineal}

La correlación estandariza la covarianza dividiéndola por el producto de las desviaciones estándar. Para una población, el coeficiente de correlación lineal es

$$
\rho_{XY} = \frac{\operatorname{Cov}(X,Y)}{\sigma_X\sigma_Y}.
$$

Para una muestra, se calcula mediante

$$
r_{XY} = \frac{s_{XY}}{s_Xs_Y},
$$

donde $s_X$ y $s_Y$ son las desviaciones estándar muestrales. Siempre que ambas variables tengan variabilidad, el coeficiente cumple $-1\leq r_{XY}\leq 1$:

\begin{itemize}
    \item $r_{XY}>0$ indica asociación lineal positiva: los valores mayores de una variable tienden a acompañarse de valores mayores de la otra.
    \item $r_{XY}<0$ indica asociación lineal negativa: los valores mayores de una variable tienden a acompañarse de valores menores de la otra.
    \item Un valor cercano a cero indica poca asociación lineal; no descarta una relación no lineal.
    \item $r_{XY}=1$ o $r_{XY}=-1$ corresponde a una relación lineal positiva o negativa perfecta en los datos.
\end{itemize}

La correlación no depende de las unidades de medida, a diferencia de la covarianza, pero describe únicamente la intensidad y dirección de la asociación lineal. Por sí sola no demuestra una relación de causa y efecto.